\documentclass{szb-solution}

\area{Analitikus geometria}
\title{Térgeometriai alakzatok}
\subject{Matematika G1}
\subjectCode{BMETE94BG01}
\date{Utoljára frissítve: \today}
\docno{2}

\begin{document}
\maketitle

\subsection{Párhuzamos egyenesek távolsága}

Az $e_1:3x-4y-10=0$ és $e_2:6x-8y+5=0$ egyenesek párhuzamosak,
mert normálvektoraik arányosak. Az $e_1$ egyenes $x=0$ koordinátájú pontja
$P=(0;-5/2)$. A pont és az $e_2$ egyenes távolsága:
$$
  d(e_1,e_2)
  =d(P,e_2)
  =\frac{|6\cdot0-8(-5/2)+5|}{\sqrt{6^2+(-8)^2}}
  =\frac{25}{10}
  =\boxed{\frac52}
  \text.
$$

\subsection{Két ponton átmenő egyenes}

A $P(-2;5;6)$ és $Q(7;-1;3)$ pontokon átmenő egyenes irányvektora
$$
  \overrightarrow{PQ}=Q-P=(9;-6;-3)=3(3;-2;-1)
  \text.
$$
Az egyenes vektoros és paraméteres alakja például
$$
  \boxed{
    \rvec r(t)=
    \begin{bmatrix}-2\\5\\6\end{bmatrix}
    +t\begin{bmatrix}3\\-2\\-1\end{bmatrix}
  },
  \qquad t\in\Reals,
$$
$$
  \boxed{
    x=-2+3t,
    \qquad y=5-2t,
    \qquad z=6-t
  }
  \text.
$$

\subsection{Metsző egyenesek paramétere}

Az egyenesek paraméteres alakjai:
$$
  \begin{aligned}
    e_1:&\quad x=-2+2t_1, & y=-3t_1, & z=1+4t_1,\\
    e_2:&\quad x=3+\alpha t_2, & y=1+4t_2, & z=7+2t_2.
  \end{aligned}
$$
A közös pontban a koordináták megegyeznek. A $y$- és $z$-koordinátákból
$$
  -3t_1=1+4t_2,
  \qquad
  1+4t_1=7+2t_2
$$
adódik, aminek megoldása $t_1=1$ és $t_2=-1$. Az $x$-koordinátára
$$
  -2+2\cdot1=3+\alpha(-1)
  \quad\Longrightarrow\quad
  \boxed{\alpha=3}
  \text.
$$

\subsection{Párhuzamos térbeli egyenesek távolsága}

Legyen
$$
  \rvec r_1=(2;-1;0),\quad \rvec v_1=(3;4;2),
  \qquad
  \rvec r_2=(7;1;3),\quad \rvec v_2=(6;8;4)=2\rvec v_1
  \text.
$$
Az egyenesek párhuzamosak, ezért
$$
  d
  =\frac{|\rvec v_2\times(\rvec r_2-\rvec r_1)|}{|\rvec v_2|}
  =\frac{
    \left|\begin{bmatrix}6\\8\\4\end{bmatrix}
    \times\begin{bmatrix}5\\2\\3\end{bmatrix}\right|
  }{\sqrt{6^2+8^2+4^2}}
  =\frac{\sqrt{16^2+2^2+(-28)^2}}{\sqrt{116}}
  =\boxed{3}
  \text.
$$

\subsection{Kitérő egyenesek normáltranszverzálisa}

Az egyenesek adatai:
$$
  \rvec r_1=(-7;4;4),\quad \rvec v_1=(3;-2;3),
  \qquad
  \rvec r_2=(1;-8;-12),\quad \rvec v_2=(1;2;-1)
  \text.
$$
A közös merőleges iránya
$$
  \rvec v_1\times\rvec v_2=(-4;6;8)=-2(2;-3;-4)
  \text.
$$
Továbbá $\rvec r_2-\rvec r_1=(8;-12;-16)=4(2;-3;-4)$, vagyis a két megadott
fixpontot összekötő szakasz már merőleges mindkét egyenesre. Így
$$
  \boxed{d=|\rvec r_2-\rvec r_1|=\sqrt{8^2+(-12)^2+(-16)^2}=4\sqrt{29}}
  \text,
$$
és a normáltranszverzális egyenlete
$$
  \boxed{
    \rvec r(u)=
    \begin{bmatrix}-7\\4\\4\end{bmatrix}
    +u\begin{bmatrix}2\\-3\\-4\end{bmatrix}
  },
  \qquad u\in\Reals
  \text.
$$

\subsection{Három pont által kifeszített sík}

Legyen $P(0;-1;2)$, $Q(2;-1;1)$ és $R(4;3;-2)$. Ekkor
$$
  \overrightarrow{PQ}=(2;0;-1),
  \qquad
  \overrightarrow{PR}=(4;4;-4),
$$
$$
  \overrightarrow{PQ}\times\overrightarrow{PR}
  =(4;4;8)\neq\nvec
  \text.
$$
A pontok tehát nem esnek egy egyenesre. A sík normálvektora lehet
$\rvec n=(1;1;2)$, ezért a $P$ ponton átmenő sík:
$$
  x+(y+1)+2(z-2)=0
  \quad\Longleftrightarrow\quad
  \boxed{x+y+2z-3=0}
  \text.
$$

\subsection{Egyenes és sík párhuzamossága}

Az egyenes irányvektora $\rvec v=(2;2;\alpha)$, a sík normálvektora pedig
$\rvec n=(1;3;-2\alpha)$. A párhuzamosság feltétele
$$
  \rvec v\cdot\rvec n
  =2+6-2\alpha^2=0
  \quad\Longrightarrow\quad
  \alpha^2=4
  \text.
$$
Tehát
$$
  \boxed{\alpha=\pm2}
  \text.
$$

\subsection{Egyenes és sík metszéspontja}

Az egyenes paraméteres alakja
$$
  x=1+t,
  \qquad y=-1-2t,
  \qquad z=6t
  \text.
$$
Behelyettesítve a $2x+3y+z-1=0$ síkegyenletbe:
$$
  2(1+t)+3(-1-2t)+6t-1=0
  \quad\Longrightarrow\quad
  2t-2=0,
  \qquad t=1
  \text.
$$
A metszéspont
$$
  \boxed{P=(2;-3;6)}
  \text.
$$

\subsection{Két sík metszésvonala}

A síkok normálvektorai $\rvec n_1=(1;-2;3)$ és
$\rvec n_2=(3;2;-5)$, ezért a metszésvonal irányvektora
$$
  \rvec v=\rvec n_1\times\rvec n_2=(4;14;8)=2(2;7;4)
  \text.
$$
A $z=0$ választással
$$
  x-2y=4,
  \qquad 3x+2y=4
  \quad\Longrightarrow\quad
  x=2,
  \qquad y=-1
  \text.
$$
Így a metszésvonal
$$
  \boxed{
    \rvec r(t)=
    \begin{bmatrix}2\\-1\\0\end{bmatrix}
    +t\begin{bmatrix}2\\7\\4\end{bmatrix}
  },
  \qquad t\in\Reals
  \text.
$$

\subsection{Három sík közös pontja}

A síkok egyenletrendszere:
$$
  \begin{cases}
    2x+y-z=2,\\
    x-3y+z=-1,\\
    x+y+z=3.
  \end{cases}
$$
Az első és a harmadik egyenlet különbségéből $x-2z=-1$ adódik. Az első két
egyenlet különbségével együtt ebből, majd visszahelyettesítéssel
$$
  \boxed{P=(1;1;1)}
  \text.
$$
A keresett sík párhuzamos az $x+y+2z=0$ síkkal, ezért normálvektora
$(1;1;2)$. A $P$ ponton átmenő egyenlete
$$
  (x-1)+(y-1)+2(z-1)=0
  \quad\Longleftrightarrow\quad
  \boxed{x+y+2z-4=0}
  \text.
$$

\end{document}
